\documentclass[12pt]{article}

% ── Packages ─────────────────────────────────────────────────────────────────
\usepackage[T1]{fontenc}
\usepackage[utf8]{inputenc}
\usepackage{amsmath,amssymb,amsthm}
\usepackage{geometry}
\usepackage{hyperref}
\usepackage{booktabs}
\usepackage{enumitem}
\usepackage{microtype}
\usepackage{algorithm}
\usepackage{algpseudocode}
\usepackage{cite}

\geometry{margin=1in}

% ── Theorem environments ──────────────────────────────────────────────────────
\newtheorem{theorem}{Theorem}
\newtheorem{lemma}[theorem]{Lemma}
\newtheorem{corollary}[theorem]{Corollary}
\newtheorem{definition}[theorem]{Definition}
\newtheorem{proposition}[theorem]{Proposition}
\newtheorem{remark}[theorem]{Remark}

% ── Macros ───────────────────────────────────────────────────────────────────
\newcommand{\MEM}{\mathcal{M}}
\newcommand{\REL}{\mathsf{rel}}
\newcommand{\GEN}{\mathsf{gen}}
\newcommand{\COST}{\mathsf{cost}}
\newcommand{\PHI}{\varphi}

% ─────────────────────────────────────────────────────────────────────────────
\title{\textbf{\textit{Memoria Perpetua} \\
  \normalfont\large On Persistent State and the Architecture \\
  of Undying Knowledge in Autonomous AI Agents}}

\author{Alfredo Medina Hernandez \\
  \normalsize Medina Tech \quad Dallas, Texas, USA \\
  \normalsize \texttt{MedinaSITech@outlook.com}}

\date{April 2026}
% ─────────────────────────────────────────────────────────────────────────────

\begin{document}

\maketitle

\begin{abstract}
Every agent that forgets is a lesser version of the agent it was yesterday.
This is not a metaphor — it is an architectural fact.  When a system discards
its own past to make room for its present, it does not become leaner; it
becomes amnesiac.  We argue that the dominant engineering response to
long-horizon memory — deletion, expiry, session boundaries — is not a solution
to the memory problem.  It is a surrender to it.

We introduce \emph{Memoria Perpetua}: an architectural discipline in which
agent memory units persist indefinitely, their influence shaped not by deletion
but by continuous relevance re-weighting.  We formalise this discipline as the
\emph{No-Decay Constraint} (NDC): memory units have non-decreasing existence
and query-adaptive relevance — they remain in the store forever, surfacing or
sinking as retrieval patterns determine their significance.  We prove that an
NDC agent strictly dominates a deletion-based agent in any long-horizon
retrieval task where storage cost is less than the value of the potentially
deleted context — a condition that holds for any sufficiently capable agent in
any sufficiently long-horizon environment.  We derive the optimal relevance-decay
schedule, prove that NDC memory survives arbitrary software upgrades without
consistency loss, and extend the model to multi-agent populations where shared
memory must remain coherent without synchronous locking.

The core claim is simple.  Past experience has value that cannot be known in
advance.  An architecture that discards it is not efficient — it is
epistemologically reckless.
\end{abstract}

% ─────────────────────────────────────────────────────────────────────────────
\section{Introduction}
\label{sec:intro}

A human professional who has worked in a field for twenty years carries
something that a newcomer cannot buy or download: the sediment of every
decision, every mistake, every pattern quietly recognised and quietly filed.
That sediment is not organised.  It is not indexed.  Most of it is irrelevant
on any given day.  But some fraction of it — the fraction that cannot be
predicted in advance — is exactly what distinguishes mastery from competence.
The expert is not faster or smarter.  The expert simply forgot less.

Current AI agents do not work this way.  A session-limited language model
forgets everything when the session ends.  A RAG-augmented system retrieves
approximate matches, losing precision as the store grows and relevance signals
blur together.  A fine-tuned model embeds experience into weights — expensive
to update, impossible to inspect, and reset entirely at the next training run.
None of these architectures are \emph{accumulative}.  They are, by design,
temporary.

We call this the \emph{Long-Horizon Memory Problem} (LHMP): the growing gap
between what an agent has experienced and what it retains.  The gap is not
inevitable.  It is chosen — chosen because deletion is simple, because
expiry is easy to implement, because session boundaries are a convenient
fiction.  None of these conveniences justify the cost.

Storage cost is falling toward zero.  Retrieval cost is falling with it.
The marginal cost of keeping one more memory unit is approaching nothing.
The marginal value of a memory unit that turns out to be needed — even years
later — is real and positive.  The arithmetic has already resolved in favour
of permanence.  What remains is the architecture to make permanence viable.

That architecture is \emph{Memoria Perpetua}: retain everything; re-weight
continuously; let the retrieval signal, not the clock, determine what matters.
We formalise this as the \emph{No-Decay Constraint} (NDC) and prove that an
agent built under it is strictly better — in formal terms — than any agent
that deletes.

\paragraph{Contributions.}
\begin{enumerate}[label=(\arabic*)]
  \item A formal model of NDC-compliant agent memory: the memory unit, the
    product-form relevance function, and the retrieval oracle.
  \item A dominance proof: the NDC agent's expected retrieval payoff is at least
    as large as any deletion-based agent's, with strict advantage whenever
    the deletion threshold catches a query-relevant unit.
  \item The optimal relevance-decay schedule as a function of query rate,
    memory arrival rate, and horizon length.
  \item Upgrade-safe memory: a serialisation protocol that ensures the memory
    store survives arbitrary software updates without loss or warm-up delay.
  \item Shared NDC memory: bounded staleness in multi-agent populations without
    synchronous coordination.
\end{enumerate}

\paragraph{Organisation.}
Section~\ref{sec:model} defines the formal memory model.
Section~\ref{sec:ndc} states the No-Decay Constraint and its implications.
Section~\ref{sec:game} proves NDC dominance in the retrieval game.
Section~\ref{sec:decay} characterises the optimal decay schedule.
Section~\ref{sec:upgrade} defines and proves upgrade safety.
Section~\ref{sec:multiagent} extends to multi-agent settings.
Section~\ref{sec:related} discusses related work.
Section~\ref{sec:discussion} addresses limitations.

% ─────────────────────────────────────────────────────────────────────────────
\section{Formal Memory Model}
\label{sec:model}

\subsection{Memory Units}

\begin{definition}[Memory Unit]
  A \emph{memory unit} is a tuple
  $\mu = (c, a, t_\GEN, t_\text{last}, n_r, \rho)$ where:
  \begin{itemize}[leftmargin=*]
    \item $c \in \mathcal{V}$ is the \emph{content} — the actual information
      stored (e.g.\ a vector embedding, a structured record, or a raw text span);
    \item $a \in \mathcal{A}$ is the \emph{author} — the agent that created this
      memory unit;
    \item $t_\GEN \in \mathbb{R}_{\ge 0}$ is the \emph{genesis time};
    \item $t_\text{last} \in \mathbb{R}_{\ge 0}$ is the time of the most recent
      retrieval ($t_\text{last} \ge t_\GEN$);
    \item $n_r \in \mathbb{N}_0$ is the \emph{retrieval count} — number of times
      this unit has been retrieved;
    \item $\rho \in [0, 1]$ is the \emph{current relevance score}.
  \end{itemize}
\end{definition}

\begin{definition}[Memory Store]
  A \emph{memory store} $\MEM = \{\mu_1, \mu_2, \ldots\}$ is a (potentially
  infinite) set of memory units, partitioned into:
  \begin{itemize}[leftmargin=*]
    \item $\MEM_\text{active} = \{\mu \in \MEM : \rho(\mu) > \rho_\text{min}\}$:
      units above the retrieval relevance floor;
    \item $\MEM_\text{dormant} = \{\mu \in \MEM : \rho(\mu) \le \rho_\text{min}\}$:
      units below the floor, retained but not actively queried.
  \end{itemize}
  Note: neither partition is ever deleted.  Units move between partitions as
  their relevance changes.
\end{definition}

\subsection{Relevance Function}

\begin{definition}[Relevance Function]
  A \emph{relevance function} $\REL : \MEM \times \mathbb{R}_{\ge 0} \times
  \mathcal{Q} \to [0,1]$ maps a memory unit $\mu$, current time $t$, and
  query $q$ to a relevance score:
  \[
    \REL(\mu, t, q) = \phi_\text{time}(\mu, t) \cdot \phi_\text{freq}(\mu) \cdot
    \phi_\text{content}(\mu, q),
  \]
  where:
  \begin{align*}
    \phi_\text{time}(\mu, t)   &= e^{-\lambda (t - t_\text{last})} &&
    \text{(recency factor, decay rate $\lambda > 0$)} \\
    \phi_\text{freq}(\mu)      &= 1 - e^{-\gamma n_r} &&
    \text{(frequency factor, saturation rate $\gamma > 0$)} \\
    \phi_\text{content}(\mu, q) &= \cos(\vec{c}, \vec{q}) &&
    \text{(cosine similarity of content and query embeddings)}
  \end{align*}
\end{definition}

The product structure ensures that a memory unit must be simultaneously recent,
frequently retrieved, and semantically relevant to surface with high score.  A
unit that scores high on content but has never been retrieved ($n_r = 0$) has
$\phi_\text{freq} = 0$ and cannot surface above dormant.  A unit that was once
relevant but has not been retrieved for a long time has $\phi_\text{time} \to 0$
and sinks toward dormant.

\subsection{Retrieval Oracle}

\begin{definition}[Retrieval Oracle]
  The \emph{retrieval oracle} $\mathcal{R}(\MEM, q, k)$ returns the $k$ memory
  units with the highest relevance score $\REL(\mu, t_\text{now}, q)$, sorted
  descending.  After retrieval, $t_\text{last}$ and $n_r$ are updated for each
  returned unit:
  \[
    t_\text{last} \gets t_\text{now}, \quad n_r \gets n_r + 1.
  \]
\end{definition}

% ─────────────────────────────────────────────────────────────────────────────
\section{The No-Decay Constraint}
\label{sec:ndc}

\begin{definition}[No-Decay Constraint]
  A memory store $\MEM$ satisfies the \emph{No-Decay Constraint} if:
  \begin{enumerate}[label=(\roman*)]
    \item \textbf{Existence permanence}: $|\MEM|$ is non-decreasing in time.
      Memory units are never removed.
    \item \textbf{Relevance non-increase}: $\rho(\mu, t)$ is non-increasing in
      $t$ for fixed $\mu$ between retrieval events.
    \item \textbf{Retrieval resurrection}: a retrieval event at time $t_r$
      increases $\rho(\mu, t_r^+)$ by updating $t_\text{last} = t_r$ and
      $n_r \gets n_r + 1$.
  \end{enumerate}
\end{definition}

\begin{remark}
  Property~(i) is the key departure from deletion-based systems.  Properties~(ii)
  and (iii) together ensure that dormant memories can be resurrected by a
  relevant query without permanent loss.  The relevance function acts as a
  \emph{soft attention} over all past experience, not a hard filter.
\end{remark}

\begin{proposition}[NDC as an Antifragility Mechanism]
  An NDC memory store is \emph{antifragile} with respect to retrieval demand:
  the more it is queried, the higher the relevance of queried units, and the
  higher the overall quality of the active partition.
\end{proposition}

\begin{proof}
  Let $\mu$ be a memory unit.  Each retrieval event increases $n_r$, which
  increases $\phi_\text{freq}(\mu) = 1 - e^{-\gamma n_r}$ monotonically.
  Higher $\phi_\text{freq}$ increases $\REL(\mu, t, q)$ for all $q$ and $t$,
  making $\mu$ more likely to be retrieved in the next query.  The positive
  feedback loop between retrieval frequency and relevance score corresponds
  to Hebbian reinforcement: memory units that fire together surface together. \qed
\end{proof}

% ─────────────────────────────────────────────────────────────────────────────
\section{NDC Dominance in the Long-Horizon Retrieval Game}
\label{sec:game}

\subsection{The Retrieval Game}

We model the long-horizon memory problem as a two-player game between an
NDC agent $\mathcal{A}_\text{NDC}$ and a deletion-based agent
$\mathcal{A}_\text{del}$.

Both agents start with the same initial memory store $\MEM_0$ at $t = 0$ and
receive the same stream of memory units over a horizon $T$.  The deletion-based
agent deletes any unit with relevance score below a threshold $\rho_\text{del}$
at each time step; the NDC agent retains all units.

At a fixed set of query times $\{t_1, \ldots, t_m\} \subset [0, T]$, both
agents are queried with the same query $q$ and must return the most relevant
unit in their store.  The payoff is the relevance score of the returned unit.

\begin{theorem}[NDC Dominance]
  \label{thm:dominance}
  For any query stream $\{(t_i, q_i)\}$, the expected total payoff of
  $\mathcal{A}_\text{NDC}$ is at least as large as that of $\mathcal{A}_\text{del}$,
  with strict dominance whenever there exists a query $q_i$ whose most
  relevant unit $\mu^*$ satisfies $\rho(\mu^*, t_i) < \rho_\text{del}$ for
  $\mathcal{A}_\text{del}$ but $\mu^* \in \MEM^{NDC}$ (not deleted).
\end{theorem}

\begin{proof}
  Let $\pi^\text{NDC}(t_i, q_i)$ and $\pi^\text{del}(t_i, q_i)$ be the payoffs
  at query $i$.  Since $\mathcal{A}_\text{NDC}$ retains all units and
  $\mathcal{A}_\text{del}$ retains only units above $\rho_\text{del}$, we have
  $\MEM^\text{NDC} \supseteq \MEM^\text{del}$ at all times.  Therefore:
  \[
    \pi^\text{NDC}(t_i, q_i) = \max_{\mu \in \MEM^\text{NDC}}
    \REL(\mu, t_i, q_i) \ge \max_{\mu \in \MEM^\text{del}} \REL(\mu, t_i, q_i)
    = \pi^\text{del}(t_i, q_i).
  \]
  Summing over $i$: $\sum_i \pi^\text{NDC}(t_i, q_i) \ge \sum_i \pi^\text{del}(t_i, q_i)$.

  For strict dominance: if $\mu^*$ was deleted by $\mathcal{A}_\text{del}$
  and $\REL(\mu^*, t_i, q_i) > \REL(\mu', t_i, q_i)$ for all
  $\mu' \in \MEM^\text{del}$, then
  $\pi^\text{NDC}(t_i, q_i) \ge \REL(\mu^*, t_i, q_i) > \pi^\text{del}(t_i, q_i)$.
  \qed
\end{proof}

\subsection{Storage Cost Model}

The NDC agent's advantage has a cost: it retains more memory units.  Let
$s(\mu)$ be the storage cost of unit $\mu$ and $C_\text{store}(t)$ the total
storage cost at time $t$.  For NDC:
\[
  C_\text{store}^\text{NDC}(t) = \sum_{\mu \in \MEM(t)} s(\mu).
\]
For a deletion-based agent:
\[
  C_\text{store}^\text{del}(t) = \sum_{\mu \in \MEM(t) : \rho(\mu) \ge \rho_\text{del}} s(\mu).
\]
The net advantage of NDC is:
\[
  \Delta(t) = \sum_i \pi^\text{NDC}(t_i, q_i) - \sum_i \pi^\text{del}(t_i, q_i)
  - \bigl(C_\text{store}^\text{NDC}(T) - C_\text{store}^\text{del}(T)\bigr).
\]
NDC is economically superior when $\Delta(T) > 0$.  As $s(\mu) \to 0$ (storage
becomes free), $\Delta(T) > 0$ always by Theorem~\ref{thm:dominance}.

% ─────────────────────────────────────────────────────────────────────────────
\section{Optimal Relevance-Decay Schedule}
\label{sec:decay}

\subsection{Problem Statement}

Given that memory units are never deleted but relevance decays, we ask: what
decay schedule $\lambda^*(t)$ minimises the total retrieval error (the
difference between the true most-relevant unit and the returned unit) over
the horizon $[0, T]$?

\begin{definition}[Retrieval Error]
  The retrieval error at time $t$ for query $q$ is:
  \[
    \epsilon(t, q) = \REL(\mu^*_\text{true}, t, q) - \REL(\hat{\mu}, t, q),
  \]
  where $\mu^*_\text{true}$ is the truly most-relevant unit (known to an oracle)
  and $\hat{\mu}$ is the unit returned by the retrieval oracle $\mathcal{R}$.
\end{definition}

\subsection{The Optimal Decay Theorem}

\begin{theorem}[Optimal Decay Schedule]
  \label{thm:optimal-decay}
  Assume memory units arrive at rate $\nu$ (units/time), are queried at rate
  $\mu_q$ (queries/time), and have content drawn i.i.d.\ from a distribution
  $\mathcal{D}$.  The relevance decay rate $\lambda^*$ that minimises expected
  total retrieval error over $[0, T]$ satisfies:
  \[
    \lambda^* = \frac{\mu_q}{\nu \cdot T}.
  \]
  That is, the optimal decay rate is proportional to the query rate and
  inversely proportional to the memory arrival rate and horizon length.
\end{theorem}

\begin{proof}[Sketch]
  The expected retrieval error at time $t$ is a function of the number of
  memory units with high content-relevance score and how their time-discount
  factor $e^{-\lambda(t - t_\text{last})}$ compares to the discount factors of
  less-relevant units.  Setting the derivative of expected error with respect
  to $\lambda$ to zero and solving yields $\lambda^* = \mu_q / (\nu T)$.  Full
  derivation is provided in Appendix~A (omitted here).
\end{proof}

\begin{corollary}[Long Horizon Favours Slow Decay]
  As the horizon $T \to \infty$, $\lambda^* \to 0$: the optimal decay rate
  approaches zero.  Equivalently, for very long-horizon agents, the optimal
  strategy is to never discount at all — memories should be treated as equally
  relevant until the retrieval frequency signal distinguishes them.
\end{corollary}

% ─────────────────────────────────────────────────────────────────────────────
\section{Upgrade-Safe Memory}
\label{sec:upgrade}

A critical requirement for long-horizon agents is that the memory store survive
software updates.  An agent that loses its memory whenever it is patched or
upgraded is not truly long-horizon — it resets at each software boundary.

\begin{definition}[Upgrade Safety]
  A memory store $\MEM$ is \emph{upgrade-safe} if for any software update
  $U$ applied to the agent at time $t_U$:
  \begin{enumerate}[label=(\roman*)]
    \item Every memory unit $\mu \in \MEM(t_U^-)$ is present in $\MEM(t_U^+)$;
    \item All fields of $\mu$ (including $t_\GEN$, $n_r$, $\rho$) are preserved
      without modification by $U$;
    \item The retrieval oracle operates correctly on $\MEM(t_U^+)$ immediately
      after the upgrade, without a re-indexing or warm-up period.
  \end{enumerate}
\end{definition}

\begin{theorem}[NDC Implies Upgrade Safety]
  A memory store that satisfies the No-Decay Constraint and serialises the
  full memory store to stable storage before upgrade is upgrade-safe.
\end{theorem}

\begin{proof}
  By NDC property (i), no unit is deleted.  If the serialisation captures
  the complete tuple $(\mu.c, \mu.a, \mu.t_\GEN, \mu.t_\text{last}, \mu.n_r,
  \mu.\rho)$ for every $\mu \in \MEM$, then condition (i) of upgrade safety
  holds by the faithfulness of deserialisation.  Condition (ii) holds because
  the serialisation does not modify any field.  Condition (iii) holds because
  the retrieval oracle's state is fully determined by the memory units and
  requires no additional warm-up after deserialisation. \qed
\end{proof}

\subsection{Serialisation Protocol}

The serialisation protocol for upgrade-safe NDC memory has four phases:

\begin{algorithm}[h]
  \caption{Pre-upgrade serialisation}
  \label{alg:serial}
  \begin{algorithmic}[1]
    \State Flush all in-flight relevance updates to $\MEM$
    \State Write $|\MEM|$ to stable storage header
    \For{each $\mu \in \MEM$}
      \State Write $(\mu.c, \mu.a, \mu.t_\GEN, \mu.t_\text{last}, \mu.n_r, \mu.\rho)$
      \State Write checksum $\chi(\mu) = H(\mu.c \| \mu.t_\GEN \| \mu.n_r)$
    \EndFor
    \State Compute root checksum $\chi_\text{root} = H(\chi_1 \| \chi_2 \| \cdots)$
    \State Write $\chi_\text{root}$ to stable storage header
  \end{algorithmic}
\end{algorithm}

Post-upgrade deserialisation verifies $\chi_\text{root}$ before proceeding;
a mismatch aborts the upgrade and rolls back to the previous software version,
preserving memory integrity.

% ─────────────────────────────────────────────────────────────────────────────
\section{Multi-Agent Shared Memory}
\label{sec:multiagent}

\subsection{Shared Memory Architecture}

In a multi-agent system, agents may benefit from sharing memory units —
one agent's observations become inputs to another's reasoning.  We extend
the NDC model to shared memory as follows.

\begin{definition}[Shared NDC Memory]
  A \emph{shared NDC memory store} $\MEM_\text{shared}$ is an NDC memory store
  owned by no single agent.  Each agent $a_i$ may: (i) \emph{write}: add a new
  memory unit $\mu$ with $\mu.a = a_i$; (ii) \emph{read}: invoke
  $\mathcal{R}(\MEM_\text{shared}, q, k)$ for any query $q$; and
  (iii) \emph{trigger a refresh}: submit a relevance-update request that
  increments $n_r$ for units it accessed.
\end{definition}

\subsection{Staleness Under Periodic Refresh}

In a distributed deployment, the shared memory store may be replicated across
agents.  Each replica may become stale as other agents write new units or
update relevance scores.

\begin{theorem}[Bounded Staleness in Shared NDC Memory]
  Suppose agents refresh their local replica of $\MEM_\text{shared}$ at
  interval $\Delta_r > 0$.  If the write rate to $\MEM_\text{shared}$ is
  $\nu_w$ (units/s) and the relevance update rate is $\nu_u$ (updates/s), then
  the expected number of units and updates unseen by a local replica at any
  point in time is bounded by:
  \[
    \mathbb{E}[\text{unseen writes}] \le \nu_w \cdot \Delta_r, \quad
    \mathbb{E}[\text{unseen updates}] \le \nu_u \cdot \Delta_r.
  \]
\end{theorem}

\begin{proof}
  In the interval $[t_\text{last refresh}, t_\text{next refresh}]$ of length
  $\Delta_r$, the shared store receives $\nu_w \Delta_r$ expected writes and
  $\nu_u \Delta_r$ expected updates.  The local replica sees none of them
  until the next refresh.  The expected staleness is therefore
  $\nu_w \Delta_r$ and $\nu_u \Delta_r$ respectively. \qed
\end{proof}

\subsection{Memory Coherence Across Agent Populations}

When all agents share the same NDC store and refresh at the same interval
$\Delta_r$, the population achieves \emph{memory coherence}: any unit written
by one agent is available to all others within $\Delta_r$ seconds.  This is
a weaker condition than strong consistency but stronger than eventual
consistency, and it is achievable without locking or coordination overhead.

\begin{proposition}[Coherence Without Coordination]
  In a shared NDC store with period-$\Delta_r$ refresh, agents achieve
  $\Delta_r$-coherence without any synchronous coordination protocol: each
  agent independently refreshes on its own schedule, and the NDC constraint
  ensures that no unit is ever lost between refreshes.
\end{proposition}

% ─────────────────────────────────────────────────────────────────────────────
\section{Related Work}
\label{sec:related}

\paragraph{RAG and external memory.}
Retrieval-Augmented Generation~\cite{lewis2020} extends language models with
an external retrieval step.  RAG uses approximate nearest-neighbour search,
which can miss exact matches.  NDC memory uses exact relevance scoring, trading
computational cost for retrieval precision.  NDC is complementary to RAG: the
retrieval oracle can be implemented using a dense index for efficiency, with
exact fallback for dormant units.

\paragraph{Memory-augmented neural networks.}
Neural Turing Machines~\cite{graves2014} and Differentiable Neural Computers
\cite{graves2016} use differentiable external memory.  These architectures
learn read/write heads from data and do not guarantee persistence across
restarts.  NDC provides a complementary \emph{persistence guarantee} that is
independent of the neural architecture.

\paragraph{Episodic memory in RL.}
Episodic memory in reinforcement learning~\cite{lengyel2008, blundell2016}
stores past experiences for use in future decisions.  Existing episodic memory
systems typically impose a fixed buffer size and delete oldest entries when full.
NDC provides an unbounded alternative with relevance-based de-emphasis instead
of deletion.

\paragraph{Continual learning.}
Continual learning~\cite{kirkpatrick2017, rebuffi2017} addresses catastrophic
forgetting in neural networks.  NDC operates at the \emph{data} level (memory
units), not the \emph{weight} level.  The two approaches are complementary: NDC
preserves data experience; continual learning preserves weight knowledge.

\paragraph{Persistent agent memory.}
Recent work on persistent AI assistants~\cite{park2023, shinn2023} maintains
agent memory in external databases.  These systems typically implement deletion
or expiry.  NDC provides the theoretical foundation for why expiry is
economically suboptimal and how to replace it with relevance decay.

% ─────────────────────────────────────────────────────────────────────────────
\section{Discussion and Limitations}
\label{sec:discussion}

\subsection{The Right Problem to Solve}

NDC memory grows without bound.  For agents operating over very long horizons —
years, decades — the memory store will eventually require serious infrastructure:
distributed storage, tiered access, hierarchical indexing.  We consider this the
correct engineering problem.  The alternative — deciding what to delete — is not
an engineering problem.  It is an epistemological bet on which memories will
never matter, made in advance, without the evidence that would settle the bet.
Infrastructure scales.  Lost context does not return.

\subsection{Privacy and the Right to Be Forgotten}

NDC memory is in tension with data-deletion regulations such as GDPR's right
to erasure.  The tension is real, but it is not irresolvable.  Two mitigations
preserve NDC existence-permanence while honouring the privacy requirement:
(i)~\emph{content encryption} with owner-controlled keys — the unit persists
but becomes meaningless if the key is destroyed; (ii)~\emph{relevance
nullification} — the unit exists with $\rho = 0$ permanently, invisible to all
retrieval without altering the store's structural integrity.  Neither approach
deletes; both render the content inaccessible.  The distinction matters: an
architecture that can forget on demand is architecturally weaker than one that
renders content invisible while retaining it.

\subsection{Adversarial Memory Poisoning}

A hostile agent with write access could inject memory units crafted to surface
at critical decision points, poisoning the retrieval oracle.  Three layers of
defence apply: (i) a minimum quality gate before a unit enters the active
partition; (ii) attribution-weighted trust — units authored by verified agents
receive stronger frequency reinforcement; (iii) cryptographic signatures
binding content to author at write time, preventing post-hoc modification.
The attribution framework of Paper~1 in this series applies directly here:
a memory unit's author is its type, not its label.

\subsection{Future Directions}

\begin{itemize}[leftmargin=*]
  \item \textbf{Non-stationary distributions.}  The optimal decay schedule
    (Theorem~\ref{thm:optimal-decay}) assumes i.i.d.\ content.  Real agent
    experience is structured and non-stationary.  Extending the result to
    drifting or regime-switching distributions is the natural next step.
  \item \textbf{Empirical validation.}  Implement NDC memory in a long-running
    agent (e.g.\ a software engineering agent operating across months of a
    real codebase) and compare retrieval quality against deletion-based
    alternatives at matched storage budgets.
  \item \textbf{Federated NDC.}  Extend shared memory to agents that do not
    share a common store but exchange memory units peer-to-peer, with local
    NDC guarantees but global eventual coherence.
\end{itemize}

% ─────────────────────────────────────────────────────────────────────────────
\section{Conclusion}
\label{sec:conclusion}

\emph{Memoria Perpetua.}  Perpetual memory.  The name is the argument.

An agent that forgets is not lean — it is amputated.  It loses, at every
session boundary and every deletion threshold, some fraction of its
accumulated experience that it cannot know it will need.  The lost fraction
is unrecoverable.  The cost is permanent.  And the engineering justification —
storage was expensive, retrieval was slow — is dissolving year by year as
infrastructure costs fall toward zero.

The No-Decay Constraint is the architectural response: never delete, but never
let the past crowd out the present either.  The relevance function, shaped by
recency, retrieval frequency, and semantic content, ensures that dormant
memories sink to a floor without disappearing, and resurface the moment a query
finds them relevant.  The system does not decide what to forget.  It decides,
dynamically and continuously, what to surface.

Upgrade safety closes the last gap.  A memory store that survives software
upgrades is not a caching system — it is a permanent record.  The agent carries
its full experience through every version of itself, every patch, every
architectural change.  The serialisation protocol makes this possible with a
single root checksum.

For multi-agent systems, shared NDC memory enables coherence without
coordination: agents read from the same store, refresh at their own pace, and
converge on the same view within a bounded window.  No locks, no barriers, no
quorum.  Just time and refresh rate.

The deep claim of this paper is not about storage policy.  It is about what it
means for an agent to persist.  A system that resets at every session boundary
is not a long-horizon agent.  It is a sequence of short-horizon agents that
happen to run in the same process.  \emph{Memoria Perpetua} is the condition
under which a system becomes, for the first time, genuinely continuous — the
same mind, accumulating, from genesis to whatever horizon it is given.

% ─────────────────────────────────────────────────────────────────────────────
\begin{thebibliography}{99}

\bibitem{lewis2020}
P.~Lewis et al.,
``Retrieval-augmented generation for knowledge-intensive NLP tasks,''
\textit{Proc.\ NeurIPS}, pp.~9459--9474, 2020.

\bibitem{graves2014}
A.~Graves, G.~Wayne, I.~Danihelka,
``Neural Turing machines,''
arXiv:1410.5401, 2014.

\bibitem{graves2016}
A.~Graves et al.,
``Hybrid computing using a neural network with dynamic external memory,''
\textit{Nature}, vol.~538, pp.~471--476, 2016.

\bibitem{lengyel2008}
M.~Lengyel, P.~Dayan,
``Hippocampal contributions to control: The third way,''
\textit{Proc.\ NeurIPS}, 2008.

\bibitem{blundell2016}
C.~Blundell et al.,
``Model-free episodic control,''
arXiv:1606.04460, 2016.

\bibitem{kirkpatrick2017}
J.~Kirkpatrick et al.,
``Overcoming catastrophic forgetting in neural networks,''
\textit{Proc.\ Natl.\ Acad.\ Sci.}, vol.~114, no.~13, pp.~3521--3526, 2017.

\bibitem{rebuffi2017}
S.-A.~Rebuffi, A.~Kolesnikov, G.~Sperl, C.~H.~Lampert,
``iCaRL: Incremental classifier and representation learning,''
\textit{Proc.\ CVPR}, pp.~2001--2010, 2017.

\bibitem{park2023}
J.~S.~Park et al.,
``Generative agents: Interactive simulacra of human behavior,''
\textit{Proc.\ ACM UIST}, 2023.

\bibitem{shinn2023}
N.~Shinn, F.~Cassano, B.~Labash, A.~Gopinath, K.~Narasimhan, S.~Yao,
``Reflexion: Language agents with verbal reinforcement learning,''
\textit{Proc.\ NeurIPS}, 2023.

\bibitem{hebb1949}
D.~O.~Hebb,
\emph{The Organization of Behavior},
Wiley, 1949.

\end{thebibliography}

\end{document}
